\section{Metodologia sperimentale}
\label{sec:metodologia}

Dopo aver analizzato teoricamente e implementato gli algoritmi di Needleman-Wunsch e Hirschberg, possiamo
procedere al loro confronto sperimentale.

L'obiettivo dell'analisi è verificare empiricamente le proprietà discusse nei capitoli precedenti,
concentrandoci sul tempo di esecuzione e sul consumo di memoria al crescere della lunghezza
delle sequenze.

Quando le risorse disponibili lo consentono, entrambi gli algoritmi vengono eseguiti sulla medesima istanza e con lo stesso modello di scoring.
Per input la cui dimensione rende impraticabile l'allocazione delle matrici complete di Needleman--Wunsch, viene eseguito soltanto Hirschberg;
tali casi vengono indicati esplicitamente nei risultati.

Nel seguito vengono descritti il dataset utilizzato, l'ambiente sul quale si è lavorato e 
il protocollo adottato per la raccolta delle misure, in modo tale da poter valutare
la validità dell'esperimento, che verrà presentato nel capitolo successivo.

\subsection{Dataset}

Per il confronto sperimentale sono state utilizzate due classi di input: sequenze sintetiche generate in modo controllato
e sequenze biologiche reali ottenute dal database NCBI Nucleotide.

Le sequenze sintetiche permettono di controllare direttamente le dimensioni $n$ e $m$ del problema e vengono quindi utilizzate
principalmente per studiare la scalabilità dei due algoritmi. La generazione degli input viene effettuata mediante un generatore
pseudocasuale deterministico, inizializzato attraverso un seed noto, così da rendere riproducibili le istanze utilizzate.
Sono state considerate tre famiglie di input sintetici:
\begin{itemize}
    \item[\texttt{random}] le due sequenze vengono generate indipendentemente sull'alfabeto $\Sigma_{\text{DNA}}=\{ \texttt{A},\texttt{C},\texttt{G},\texttt{T} \}$;
    \item[\texttt{identiche}] la seconda sequenza coincide con la prima;
    \item[\texttt{mutate}] la seconda sequenza viene ottenuta dalla prima sostituendo ciascun simbolo con probabilità $0.05$. In caso di
        sostituzione viene scelta una delle altre tre basi.
\end{itemize}

Sono state considerate sequenze di lunghezza crescente per generare input quadrati
\[
    250,\;500,\;1000,\;2000,\;4000,\;8000,\;12000.
\]
Per ogni combinazione di dimensione e famiglia vengono generate sette istanze indipendenti per configurazione, ottenute mediante seed differenti.
Per ciascun seed entrambi gli algoritmi vengono quindi eseguiti esattamente sulla stessa coppia di sequenze. 

Sono inoltre considerate alcune istanze rettangolari, con $n \neq m$, al fine di osservare il comportamento degli algoritmi anche quando le due
sequenze hanno lunghezze significativamente differenti.

Agli input sintetici sono affiancate quattro coppie di sequenze biologiche reali, scelte in modo da coprire dimensioni
progressivamente maggiori.

\begin{table}[htbp]
    \centering
    \small
    \begin{tabularx}{\textwidth}{
        >{\raggedright\arraybackslash}p{0.21\textwidth}
        >{\raggedright\arraybackslash}X
        >{\raggedright\arraybackslash}X
        >{\centering\arraybackslash}p{0.13\textwidth}
    }
        \toprule
        \textbf{Dataset} &
        \textbf{Sequenza 1} &
        \textbf{Sequenza 2} &
        \textbf{Lunghezza} \\
        \midrule

        RNA ribosomiale batterico &
        \textit{E. coli} \newline
        \texttt{NR\_112059.1} &
        \textit{S. enterica} \newline
        \texttt{NR\_116005.1} &
        $\sim 1.5 \cdot 10^{3}$ \\

        \addlinespace

        HIV-1 &
        HXB2 (subtipo B) \newline
        \texttt{NC\_001802.1} &
        ETH2220 (subtipo C) \newline
        \texttt{U46016.1} &
        $\sim 10^{4}$ \\

        \addlinespace

        SARS-CoV-2 &
        Wuhan-Hu-1 \newline
        \texttt{NC\_045512.2} &
        Delta (B.1.617.2) \newline
        \texttt{OK091006.1} &
        $\sim 3 \cdot 10^{4}$ \\

        \addlinespace

        Genomi plastidici &
        \textit{A. thaliana} \newline
        \texttt{NC\_000932.1} &
        \textit{O. sativa} \newline
        \texttt{NC\_001320.1} &
        $\sim 1.5 \cdot 10^{5}$ \\

        \bottomrule
    \end{tabularx}
    \caption{Coppie di sequenze biologiche reali utilizzate nel confronto sperimentale.}
    \label{tab:dataset-reali}
\end{table}

Le sequenze sono state scaricate manualmente in formato FASTA da NCBI Nucleotide utilizzando gli accession number indicati.

Nella prima coppia sono presenti anche dei simboli \texttt{N}, mentre nella sequenza Delta (B.1.617.2) è presente il simbolo \texttt{Y}, che indicano nucleotidi non determinati, questi vengono sostituiti preliminarmente 
con simboli scelti pseudocasualmente da $\Sigma_{\text{DNA}}$. Il preprocessing utilizza un seed fissato e viene eseguito una sola
volta prima della raccolta delle misure, così da garantire che entrambi gli algoritmi operino sulle medesime sequenze e che l'esperimento sia riproducibile.
Le sequenze reali sono utilizzate come carichi di lavoro realistici di dimensione crescente. Lo score ottenuto non viene pertanto 
interpretato dal punto di vista biologico; l'obiettivo del loro utilizzo è valutare sperimentalmente tempo di esecuzione e consumo di memoria
dei due algoritmi.

\subsection{Ambiente hardware e software}

Tutti gli esperimenti sono stati eseguiti sul medesimo sistema, in modo da mantenere costanti le principali condizioni che possono
influenzare le prestazioni.

\begin{table}[h]
    \centering
    \begin{tabular}{ll}
        \hline
        \textbf{Componente} & \textbf{Configurazione} \\
        \hline
        Processore & Snapdragon(R) X 12-core X1E80100 @ 3.40 GHz \\
        Architettura & ARM64 (AArch64) \\
        Memoria RAM & 16,0 GB \\
        Memoria visibile a WSL 2 & 9.7 GiB \\
        Swap WSL 2 & 4.0 GiB \\
        Sistema operativo host & Microsoft Windows 11 \\
        Ambiente di esecuzione & WSL 2 (2.6.1.0) -- Ubuntu 24.04.5 LTS \\
        Compilatore & GCC \\
        Versione compilatore & 13.3.0 \\
        Flag di compilazione & \texttt{-O2 -Wall -Wextra -Wpedantic -std=c99} \\
        \hline
    \end{tabular}
    \caption{Ambiente utilizzato per gli esperimenti.}
    \label{tab:ambiente}
\end{table}

Sebbene il sistema host disponga di 16 GB di memoria fisica, l'ambiente WSL 2 utilizzato per gli esperimenti dispone di circa 9.7 GiB di memoria RAM,
oltre a 4.0 GiB di spazio di swap. Tali valori costituiscono quindi il riferimento effettivo per la valutazione della fattibilità delle esecuzioni.

\subsection{Misura dei tempi e della memoria}
\label{sec:misura}

Ogni esecuzione del benchmark avviene in un processo separato.
La misura temporale comprende esclusivamente l'esecuzione dell'algoritmo di allineamento. La generazione di input sintetici, la lettura
e il parsing dei file FASTA e le successive verifiche di correttezza vengono pertanto effettuati al di fuori della regione temporizzata.

Vengono raccolte due misure temporali. Il tempo reale trascorso, o \emph{wall time}, viene misurato utilizzando \texttt{CLOCK\_MONOTONIC},
mentre il tempo di CPU viene ottenuto mediante \texttt{CLOCK\_PROCESS\_CPUTIME\_ID}. Il wall time viene utilizzato come misura principale delle prestazioni
temporali, mentre il CPU time costituisce una misura complementare.

Per gli input sintetici, l'ordine delle esecuzioni di Needleman--Wunsch e Hirschberg viene alternato tra istanze successive, così da evitare
di favorire sistematicamente uno dei due algoritmi a causa dell'ordine di esecuzione.

Per le sequenze biologiche reali la coppia di input rimane invece invariata tra le diverse esecuzioni. Anche in questo caso viene alternato
l'ordine dei due algoritmi quando entrambi possono essere eseguiti.

Il consumo di memoria viene misurato tramite il massimo \emph{resident set size} del processo, ottenuto mediante \texttt{/usr/bin/time}.
La misura, espressa in KiB, rappresenta il massimo quantitativo di memoria residente raggiunto durante l'intera esecuzione del processo
e comprende quindi, oltre alle strutture ausiliarie dell'algoritmo, anche input, output, runtime e altro overhead del processo.
Essa non coincide pertanto con lo spazio ausiliario teorico analizzato nei capitoli precedenti, ma consente di confrontare empiricamente
la diversa crescita della memoria richiesta dalle due implementazioni eseguite nelle medesime condizioni.

Al termine viene verificata la validità strutturale dell'allineamento prodotto. Eliminando i gap dalle due sequenze allineate
devono essere ricostruiti esattamente gli input originali; inoltre lo score dell'allineamento viene ricalcolato e confrontato con quello
restituito dall'algoritmo.
Quando una stessa istanza viene elaborata da entrambi gli algoritmi, viene verificato inoltre
che Needleman--Wunsch e Hirschberg restituiscano lo stesso score ottimo.

Le esecuzioni terminate correttamente vengono aggregate per configurazione. Per wall time, CPU time e memoria vengono calcolate la mediana,
media e deviazione standard. La mediana viene usata come indicatore principale, poiché meno sensibile a eventuali esecuzioni anomale. Le esecuzioni
fallite o gli eventuali esaurimenti di memoria vengono invece registrati separatamente e non vengono inclusi nelle statistiche delle esecuzioni valide.
